% Technical Report -- Multi-Client File Transfer
%
% Compile on Overleaf (overleaf.com, free, no install): create a new
% project, paste this file's content as main.tex, click Recompile.
% Or locally with a TeX distribution: pdflatex technical_report.tex
%
% NOTE: Sections 1 (Introduction), 3 (Implementation), 7 (Challenges) and
% 8 (Conclusion) contain a clearly marked TODO placeholder instead of
% real prose. Per the assignment's AI policy, and because you need to be
% able to explain this report yourself, those need to be written by you,
% in your own words, before this is submitted.

\documentclass[11pt]{article}

\usepackage[margin=1in]{geometry}
\usepackage{longtable}
\usepackage{array}
\usepackage{booktabs}
\usepackage{enumitem}
\usepackage{hyperref}
\usepackage{parskip}
\usepackage{titlesec}
\usepackage{xcolor}
\usepackage{framed}

\hypersetup{colorlinks=true, linkcolor=blue, urlcolor=blue}
\titleformat{\section}{\large\bfseries}{\thesection}{1em}{}
\titleformat{\subsection}{\normalsize\bfseries}{\thesubsection}{1em}{}
\definecolor{shadecolor}{gray}{0.92}

\newcommand{\cmd}[1]{\texttt{#1}}
\newenvironment{todobox}
  {\begin{shaded}\textbf{TODO -- write this yourself:} }
  {\end{shaded}}

\title{Technical Report\\
\large Multi-Client File Transfer}
\author{}
\date{}

\begin{document}
\maketitle

\tableofcontents

\section{Introduction}

\begin{todobox}
What this application does, at a high level: a TCP server and CLI
client supporting login, private/broadcast messaging, server-stored
file upload/download, and direct client-to-client file transfer
relayed through the server, with 5+ concurrent clients supported. Write
this in your own words.
\end{todobox}

\section{Design}

\subsection{Architecture}

\begin{verbatim}
                 +--------------------+
                 |       Server       |
                 |  (server/server.py)|
                 |                    |
                 |  accept() loop     |
                 |  1 thread / client |
                 |  active_clients{}  |  <- username -> socket, guarded
                 |                    |     by clients_lock
                 |  server/storage/   |  <- UPLOAD/DOWNLOAD file store
                 +----+----------+----+
                      |          |
        TCP           |          |          TCP
   +------------------+          +------------------+
   |                                                 |
+--v-------------------+                +------------v----------+
|   Client (Alice)      |                |   Client (Bob)        |
| client/client.py      |                | client/client.py      |
|                        |                |                        |
| reader thread: owns    |                | reader thread: owns    |
|   all socket reads;    |                |   all socket reads     |
|   prints MSG/BROADCAST |                |                        |
|   pushes, queues       |                |                        |
|   replies              |                |                        |
|                        |                |                        |
| main thread: input()   |                | main thread: input()   |
|   loop, owns all       |                |   loop                 |
|   socket writes        |                |                        |
+------------------------+                +------------------------+
\end{verbatim}

See the Protocol Specification document for the full wire-format
specification.

\subsection{Key Design Decisions}

\begin{itemize}[leftmargin=1.5em]
  \item \textbf{One thread per client} on the server, coordinated
        through a single \cmd{clients\_lock}-guarded dict
        (\cmd{active\_clients}) mapping username $\to$ socket. This is
        what lets \cmd{MSG}, \cmd{BROADCAST}, and \cmd{SEND} look up and
        write to another client's live socket from a different thread.
  \item \textbf{Client read/write split}: a client's background
        \textit{reader} thread owns every \cmd{recv()} on the socket
        (including pulling file bytes off the wire for \cmd{DOWNLOAD}
        and \cmd{SEND\_INCOMING}), while the main thread owns every
        \cmd{send()}. Because TCP's two directions are independent,
        this needs no locking around the socket itself, and it lets the
        client display an incoming \cmd{MSG}/\cmd{BROADCAST}/file push
        at any time, even while the user is mid-typed-command.
  \item \textbf{Length-prefixed binary transfer}: control messages are
        line-delimited text, but file bytes are raw with an exact byte
        count announced beforehand --- avoids ambiguity between binary
        data and the \cmd{|}/\cmd{\textbackslash{}n} delimiters used by
        the text protocol.
  \item \textbf{\cmd{SEND} is a live relay, not store-and-forward}: the
        server pulls the file off the sender's socket and immediately
        pushes it to the target's socket without touching disk.
        Simpler and faster, at the cost of requiring the target to be
        online at send time.
\end{itemize}

\section{Implementation}

\begin{todobox}
Describe the socket programming techniques you'd highlight, in your own
words -- e.g.\ \texttt{SO\_REUSEADDR}, blocking sockets + threading (vs.\
\texttt{select}/\texttt{asyncio}), \texttt{sendall} vs.\ \texttt{send},
the line-reader implemented byte-by-byte, thread safety around
\texttt{active\_clients}.
\end{todobox}

\subsection{Source Layout}

\begin{longtable}{@{}p{5.2cm}p{9.3cm}@{}}
\toprule
\textbf{Path} & \textbf{Contents} \\
\midrule
\endhead
\cmd{server/server.py} & Server: connection accept loop, per-client thread, command dispatch \\
\cmd{server/protocol.py} & Wire-format helpers (identical copy on the client side) \\
\cmd{client/client.py} & Interactive CLI client: reader thread + command loop \\
\cmd{client/protocol.py} & Wire-format helpers \\
\cmd{docs/protocol\_spec.tex} & Full protocol specification \\
\cmd{experiments/test1\_rtt.py} & Performance Test 1 (RTT) \\
\cmd{experiments/test2\_broadcast\_latency.py} & Performance Test 2 (broadcast latency to 5 clients) \\
\cmd{experiments/test3\_throughput.py} & Performance Test 3 (concurrent download throughput) \\
\cmd{tests/test\_protocol.py} & End-to-end functional test covering every command and error path \\
\bottomrule
\end{longtable}

\section{Experimental Results}

All measurements below are from a real run against a live server on
localhost (\cmd{127.0.0.1:5050}), one client machine acting as both
client and server (loopback). A real multi-machine network would add
real propagation delay and NIC-level bandwidth limits that loopback
doesn't have, so these numbers represent a best case, not a
real-network prediction.

\subsection{Test 1 --- RTT (1 client, request/response, 100 repetitions)}

\begin{verbatim}
Test 1: RTT over 100 requests
  min:    0.042 ms
  avg:    0.054 ms
  median: 0.047 ms
  max:    0.212 ms
\end{verbatim}

Sub-millisecond RTT is expected on loopback --- there's no physical
network hop, so the time measured here is essentially thread scheduling
and socket syscall overhead, not network latency. The \texttt{max}
(0.212~ms) being roughly 4$\times$ the \texttt{min} is consistent with
occasional OS thread-scheduling jitter rather than anything
protocol-related.

\subsection{Test 2 --- Broadcast Latency (server $\to$ 5 clients)}

\begin{verbatim}
Test 2: broadcast delivery latency to 5 clients
  avg: 0.685 ms
  min: 0.543 ms
  max: 0.821 ms
\end{verbatim}

Latency here is higher than the single-request RTT in Test~1, which
makes sense: a \cmd{BROADCAST} requires the server to iterate over 5
sockets and call \cmd{send\_line} on each one sequentially (see
\cmd{handle\_broadcast} in \cmd{server/server.py}) before any of the 5
receiver threads can even start reading their copy, so the last
receiver's timestamp includes the fan-out delay for every client
relayed to before it.

\subsection{Test 3 --- Throughput (10 clients downloading concurrently)}

\begin{verbatim}
1 MB file:
  total transferred: 10.00 MB
  elapsed:           0.035 s
  server throughput: 286.81 MB/s

100 MB file:
  total transferred: 1000.00 MB
  elapsed:           4.107 s
  server throughput: 243.46 MB/s

1000 MB (1 GB) file:
  total transferred: 10000.00 MB
  elapsed:           221.621 s
  server throughput: 45.12 MB/s
\end{verbatim}

Throughput drops sharply as file size grows --- from $\sim$287~MB/s at
1~MB to just $\sim$45~MB/s at 1~GB, roughly a 6$\times$ fall-off, not
the flat or gently declining curve you might expect from network
overhead alone. The likely cause is in \cmd{handle\_download}
(\cmd{server/server.py}): the server reads the \textbf{entire} file
into memory (\cmd{data = f.read()}) before sending it in one
\cmd{send\_exact} call. With 10 clients concurrently downloading a
1~GB file, that's up to 10~GB held in memory across 10 threads
simultaneously, which pushes the machine into heavy memory
pressure/paging well before the network itself becomes the bottleneck.
At 1~MB and 100~MB per client (10~MB and 1000~MB total resident at
once), memory pressure is much less of a factor, which is consistent
with those two runs both staying above 240~MB/s. A straightforward fix
would be to stream the file in fixed-size chunks (e.g.\ 1~MB at a time)
instead of loading it whole, bounding memory use regardless of file
size or client count --- noted here as a known limitation rather than
implemented, since it wasn't required for this assignment's scope.

\section{Prompts of AI Use}

AI (Claude, via Claude Code) was used to generate individual pieces of
code in small increments, per the assignment's AI policy. Every prompt
used is listed below, in the order the code was produced. Each one
produced a distinct, reviewed chunk that was tested before moving to the
next.

\begin{enumerate}[leftmargin=1.7em]

\item \textit{``Write a small Python module with wire-format helpers
for a TCP app: send/receive a \texttt{COMMAND|FIELD|FIELD} text line
terminated by \texttt{\textbackslash{}n}, plus helpers to send/receive
an exact number of raw bytes for file payloads.''} $\to$
\cmd{server/protocol.py} / \cmd{client/protocol.py}

\item \textit{``Write a multi-threaded TCP server in Python. It should
accept connections in an infinite loop, spawn a thread per client, keep
a thread-safe registry of active username$\to$socket, and implement
LOGIN (check against a simple username:password text file,
auto-register unknown users, reject already-logged-in usernames),
LOGOUT, and LIST (return comma-separated active usernames).''} $\to$
\cmd{server/server.py} (skeleton)

\item \textit{``Add MSG \textless{}user\textgreater{}
\textless{}text\textgreater{} (private message forwarded to a specific
logged-in user, with an error if that user isn't online) and BROADCAST
\textless{}text\textgreater{} (forwarded to every other logged-in user)
to the server's command dispatch.''} $\to$ \cmd{server/server.py}
(messaging)

\item \textit{``Add UPLOAD \textless{}filename\textgreater{}
\textless{}filesize\textgreater{} (client pushes a file to server
storage: server replies READY, reads exactly filesize bytes, saves it,
replies UPLOAD\_OK) and DOWNLOAD \textless{}filename\textgreater{}
(server replies DOWNLOAD\_OK with the size then streams the bytes, or
ERROR FILE\_NOT\_FOUND) to the server.''} $\to$ \cmd{server/server.py}
(file transfer)

\item \textit{``Add SEND \textless{}targetuser\textgreater{}
\textless{}filename\textgreater{} \textless{}filesize\textgreater{} to
the server: relay a file live from the sending client straight to the
target client's socket if the target is online (READY, read bytes,
push SEND\_INCOMING header + bytes to target, confirm SEND\_OK to
sender); if the target isn't online, reply ERROR USER\_NOT\_FOUND.''}
$\to$ \cmd{server/server.py} (client-to-client relay)

\item \textit{``Write a TCP client with a background reader thread
that owns all socket reads: it prints incoming MSG/BROADCAST pushes
immediately, and puts any other line (a reply to a command the main
thread just sent) onto a queue. The main thread has an interactive
loop parsing user commands (login, logout, list, msg, broadcast) and
blocks on the queue for each command's reply.''} $\to$
\cmd{client/client.py} (skeleton)

\item \textit{``Add upload and download commands to the client.
Upload: send UPLOAD header, wait for READY on the reply queue, then
send the file bytes and wait for UPLOAD\_OK. Download: send DOWNLOAD
header; since the reader thread also has to pull the raw file bytes
off the socket for a DOWNLOAD\_OK reply, have the reader thread do
that read-and-save itself before queuing the reply.''} $\to$
\cmd{client/client.py} (file transfer)

\item \textit{``Add a send command to the client (send SEND header
with target/filename/size, wait for READY, then push the file bytes
and wait for SEND\_OK), and make the reader thread's SEND\_INCOMING
handler actually read the pushed file's bytes and save it to a
received/ folder instead of just printing a placeholder.''} $\to$
\cmd{client/client.py} (client-to-client send)

\item \textit{``Write three standalone scripts against the running
server: (1) one client repeatedly round-trips a LIST request and
measures RTT (min/avg/max), (2) one sender broadcasts a timestamped
message to 5 already-connected clients and each measures its own
receive latency, averaged, (3) N clients concurrently DOWNLOAD a file
of a given size and the script measures aggregate server throughput in
MB/s.''} $\to$ \cmd{experiments/test1\_rtt.py},
\cmd{experiments/test2\_broadcast\_latency.py},
\cmd{experiments/test3\_throughput.py}

\item \textit{``Write a 2--3 page protocol specification document
(message format, full command/response table, error codes, and
sequence diagrams for login, private message, upload/download, and
client-to-client send) describing the COMMAND|FIELD|... protocol I
just implemented.''} $\to$ \cmd{docs/protocol\_spec.tex}

\item (Bug investigation, not a generation prompt) Diagnosed a reported
``disconnected from server'' message appearing after a normal
\cmd{logout}. Root cause: the client's reader thread noticed the
server closing its socket right after \cmd{LOGOUT\_OK} and
misreported it as an unexpected drop. Fixed with an
\cmd{expected\_close} flag, set \textit{before} \cmd{LOGOUT} is sent
(not after the reply, which raced against the server's close). Also
added a try/except around the server's per-command dispatch so any
future unexpected exception prints a full traceback to the server
console and reports \cmd{ERROR|SERVER\_ERROR} to the client, instead
of silently dropping the connection. $\to$ \cmd{client/client.py},
\cmd{server/server.py}

\item \textit{``Add a FILES command to the server (returns
comma-separated list of filenames currently in server storage) and a
matching \texttt{files} command to the client that displays them.''}
$\to$ \cmd{server/server.py}, \cmd{client/client.py}

\item \textit{``Add an ALLUSERS command to the server that returns
every registered account (from users.txt) along with whether each one
is currently online or offline, so it's a proper directory, not just
currently-connected users like LIST/FILES.''} $\to$
\cmd{server/server.py}, \cmd{client/client.py}

\end{enumerate}

Each chunk was compiled and run against a live server before moving to
the next (see Section~6 for the concrete verification done). All code
was reviewed and is understood well enough to be explained
command-by-command in the demo.

\section{Verification Performed}

\begin{itemize}[leftmargin=1.5em]
  \item \cmd{python -m py\_compile} on all four source files --- no
        syntax errors.
  \item \cmd{tests/test\_protocol.py}: an end-to-end functional test
        that starts a real server subprocess and drives two real
        client connections through every command and every documented
        error path (login, duplicate username, wrong password, list,
        allusers, private message, broadcast, upload, download,
        download of a missing file, files listing, client-to-client
        send, send to an offline user, logout) --- \textbf{25/25
        checks passed}.
  \item All three performance scripts (\cmd{experiments/test*.py})
        were run end-to-end against a live server and produced real
        measurements (not just ``it doesn't crash'').
\end{itemize}

\begin{todobox}
After your own demo run with 5+ concurrent interactive clients, note
anything you observed (e.g.\ did LOGOUT during an active BROADCAST
work cleanly? any timing issues?).
\end{todobox}

\section{Challenges}

\begin{todobox}
Problems you personally ran into and how you solved them -- e.g.\
reasoning about which thread owns socket reads vs.\ writes, deciding on
the \texttt{SEND\_INCOMING} push design, testing concurrency, etc. This
section should reflect your own understanding, since you need to be
able to explain it live.
\end{todobox}

\section{Conclusion}

\begin{todobox}
What you learned about application-layer protocol design, TCP socket
programming, and concurrency from this assignment.
\end{todobox}

\end{document}
